The recipient-selection comparison has a structural limit. For valid original artifacts in this scope-complete graph, changed-contract roles are exactly the roles invalidated by the global check. Both methods then use the same topological order, packet and fallback. The compatible-parser audit confirms identical request sequences, so sampled differences cannot establish a distinctive routing algorithm. A redundant extra inspection in the global implementation is not a fundamental complexity benefit.

In the eight-case exception-corruption stress, completions are 6/8 for global state, 7/8 for broadcast, targeted and sparse, and 8/8 without detailed issue feedback. Targeted and sparse use two rounds; broadcast averages 1.375. This unfavorable control does not support a benefit from more corrective detail. Ordinary-condition broadcast rewrites 35 unaffected artifacts, whereas targeted, sparse and global checking rewrite none. Restoration of a deliberately corrupted appendix is necessary and is not included in that interpretation.

The primary evidence contains 152 groups of repeated identical request payloads. Eleven vary in parsed text and one varies in tool-envelope shape. The April pair has equal normalized tool parameters but different strict-parser acceptance. A May pair instead returns missing arguments versus an explicit usable tool envelope, so one normalized representation still differs. Seeds therefore do not promise bitwise reproduction. In addition, four accepted structured reports for the same project call stock codes customers. The numerical guard does not validate all prose. This directly limits any claim of full report correctness.

The real transactions ground the computations, not the authored workflow or user behavior. Scope and dependencies are registered explicitly. The parameterized pipeline is perfect without LLM calls. We therefore establish a useful execution boundary and an inspectable application, rather than a demonstrated need for multiple LLM agents or measured human attention savings.
